% !PRINTED 284
\weberpage{284}

Dann ist die gesammte Gruppe $L_7$ so dargestellt:

\weq{9}{\tau^\varrho \chi^\lambda \omega^\mu \Theta^\nu, \quad \begin{aligned} \varrho &= 0, 1, 2, 3, 4, 5. 6\\ \lambda &= 0, 1, 2; \ \mu = 0, 1\\ \nu &= 0, 1, 2, 3. \end{aligned}} % [sic] full stop printed between 5 and 6

Dass alle diese Elemente (9) von einander verschieden sind, ergiebt sich einfach daraus, dass keine Potenz von $\tau$, deren Exponent nicht durch 7 theilbar ist, in der Octa\"edergruppe enthalten sein kann, weil die Octa\"edergruppe kein Element vom $7^{\text{ten}}$ Grade hat.

Die in der Gruppe $L_7$ enthaltene Gruppe (5) vom $21^{\text{sten}}$ Grade ist dann in der Form $\tau^\varrho \chi^\lambda$ dargestellt.

Um die Composition der Elemente (9) zu charakterisiren, genügt es offenbar, wenn für jedes $\varrho$ die Elemente

\weq{10}{\Theta\tau^\varrho, \ \omega\tau^\varrho, \ \chi\tau^\varrho}

in der Form (9) dargestellt sind, weil dadurch, zusammen mit den Compositionen in der Octa\"edergruppe, das symbolische Product je zweier Elemente der Form (9) wieder in der Form (9) dargestellt werden kann.

Da nun

\weq{11}{\tau^\varrho = \begin{pmatrix} 1, & \varrho \\ 0, & 1 \end{pmatrix}.}

ist, so erhält man zunächst sehr einfach aus (6)

\weq{12}{\chi\tau^\varrho = \tau^{4\varrho}\chi,}

und daraus

\weq{13}{\chi^2\tau^\varrho = \tau^{2\varrho}\chi^2,}

und man sieht leicht, dass diese sechs Relationen eine Folge von der einen sind:

\weq{14}{\chi\tau = \tau^4\chi.}

Die übrigen Compositionen (10) erhält man aber nur durch wirkliche Ausrechnung in den einzelnen Fällen, wobei die Tabelle (7) gute Dienste leistet. Man wendet sie in der Weise an, dass man für jeden Werth $\varrho$ die Elemente

\[
\tau^{\varrho'} \omega \tau^\varrho, \quad \tau^{\varrho'} \Theta \tau^\varrho
\]

bildet, und $\varrho'$ so bestimmt, dass sich diese Elemente in der Tabelle finden, was immer nur auf eine Art möglich ist.
